
\documentclass[12pt]{article}

%%% These are some packages that are useful
\usepackage{amsfonts, lipsum}
\usepackage[scr]{rsfso}
\usepackage{amsmath,amssymb, amscd,amsbsy, bbm, amsthm, enumerate}
\usepackage[top=1in, bottom=1in, left=.45in, right=.45in]{geometry}
\usepackage[unicode]{hyperref}
\usepackage{tikz, xcolor, fancyhdr}
\usepackage{graphicx, caption, subcaption}

%%% Page formatting
%\setlength{\headsep}{30pt}
\setlength{\parindent}{0pt}
\setlength{\textheight}{9in}



%%% These define our question environment and help number things correctly
\theoremstyle{definition}
\newtheorem{thm}{Theorem}
\newtheorem{question}[thm]{Question}
\newtheorem{questionS}[thm]{***Question}
\newtheorem{prop}[thm]{Proposition}
\newtheorem{lem}[thm]{Lemma}
\newtheorem{cor}[thm]{Corollary}
\newtheorem{defn}[thm]{Definition}
\newtheorem{rem}[thm]{Remark}
\newtheorem{ex}[thm]{Example}

%%% These are some shortcuts that are handy
\def\real{{\mathbb R}}
\def\Natural{\mathbb{N}}
\def\dx{\textnormal{dx}}
\def\dy{\textnormal{dy}}
\def\dz{\textnormal{dz}}
\def\dt{\textnormal{dt}}
\def\Re{\textnormal{Re}}
\def\Im{\textnormal{Im}}
\def\exp{\textnormal{exp}}
\def\interior{\textnormal{interior}}
\def\al{\alpha}
\def\del{\delta}
\def\Del{\Delta}
\def\gam{\gamma}
\def\Gam{\Gamma}
\def\Om{\Omega}
\def\ep{\varepsilon}
\def\lam{\lambda}
\def\rational{{\mathbb Q}}
\def\integer{{\mathbb Z}}
\def\Q{{\mathbb Q}}
\def\Z{{\mathbb Z}}
\def\N{{\mathbb N}}
\def\R{{\mathbb R}}
\def\grad{\nabla}
\def\C{\mathcal C}
\def\P{\mathbb{P}}
\def\T{\mathcal T}
\def\I{\mathcal I}
\newcommand{\abs}[1]{\left| #1 \right|}
\newcommand{\inner}[1]{\langle #1 \rangle}
\newcommand{\norm}[1]{\left\lVert#1\right\rVert}
\newcommand{\spanvect}{\textnormal{span}}
\newcommand{\union}{\cup}
\newcommand{\Union}{\bigcup}
\def\intersect{\cap}
\def\Intersect{\bigcap}
\renewcommand{\neg}{\mathord{\sim}}
\newcommand{\ds}{\displaystyle}





%%%%%%%%%%%%%%%%%%%%%%%%%%%%%%%%%
%%%%%%%%%%%%%%%%%%%%%%%%%%%%%%%%%
\newcommand{\name}{Duff Baker-Jarvis}
%%%%%%%%%%%%%%%%%%%%%%%%%%%%%%%%%
%%%%%%%%%%%%%%%%%%%%%%%%%%%%%%%%%
%%% Document Starts now
\begin{document}
\title{Probability and Statistics}
\date{\today}
\author{Dr. Duff Baker-Jarvis}
\maketitle


Exercise: Given a uniform distribution $ U(a,b) $ where $ a $ is the center and $ b $ is the total width, how many data points does one need to have a 95\% confidence interval for the average have a radius of less than 5\% of the width?

Let $ S $ be a sample of $ n $ datapoints sampled randomly from $ U(a,b) $. For two points, what is the distribution for the sum of two?
A single distribution has the formula:
$$ f_U(x)= \begin{cases}
    0, & x<0 \text{ or } x>1\\
1, & 0\leq x \leq 1
\end{cases}$$

Recall the following:

\begin{thm}
The pdf for the distribution of $ X+Y $ of two random variables is given by
$$ f_{X+Y}(x)=\int_{-\infty}^{\infty} f_X(t)f_y(x-t)  dt=f_X * f_Y $$ 
\end{thm}

Thus we get
$$f_{2U}(x)= \int_{-\infty}^{\infty} 1_{[0,1]}(t)1_{[0,1]}(x-t)  dt  $$ 

Now note that $$ 1_{[0,1]}(t)1_{[0,1]}(x-t) = \begin{cases}
    0	&	\text{ if } t\not \in [0,1] \text{ or } x-t\not \in [0,1]\\
    1	&	\text{ if } t, x-t \in [0,1]
\end{cases} = 1_{[x-1,x]\cap[0,1]}
 $$ 




\end{document}
